% ============================================================
%  Chapter 5 — Factory Robots
%  Robot World: Meet the Machines That Help Us
% ============================================================

\chapter{Factory Robots}

\chapteropening{%
  Walk into a car factory and you will see hundreds of robot arms
  welding, painting, and bolting \textemdash{} all day, all night, never stopping
  for a snack.
}
\chapterbadge{Bolt's Home: The Assembly Line}

% --- Introduction ---
\section*{The Robot Factory}

Factories were the first places robots found real work. In the 1960s,
engineers realised that robots could do repetitive, dangerous, or
extremely precise jobs better than humans.

\character{Bolt}{This is where I live! I weld 400 metal joints every hour.
  My human teammates handle the creative parts \textemdash{}
  I handle the heavy lifting.}

\sectionrule

% --- What factory robots do ---
\section*{Jobs on the Line}

\begin{itemize}
  \item \term{Welding} \textemdash{} joining metal pieces with intense heat.
  \item \term{Painting} \textemdash{} spraying perfectly even coats of paint.
  \item \term{Assembly} \textemdash{} fitting parts together (screws, clips, glue).
  \item \term{Pick and place} \textemdash{} grabbing items from a conveyor belt and placing them in boxes.
  \item \term{Quality inspection} \textemdash{} using cameras to spot defects smaller than a grain of sand.
  \item \term{Palletising} \textemdash{} stacking boxes onto pallets ready for shipping.
\end{itemize}

\begin{funfact}
  A single car factory can have over 1,000 robots. Together they can
  build a complete car body in under one minute!
\end{funfact}

\sectionrule

% --- Cobots ---
\section*{Cobots \textemdash{} Robots That Work With People}

Older factory robots work behind safety cages because they are so
strong they could hurt someone by accident. But a new generation of
\term{collaborative robots} (cobots) are designed to work \emph{beside}
humans safely.

Cobots have:
\begin{itemize}
  \item Padded, rounded bodies (no sharp edges).
  \item Force sensors that stop them instantly if they bump someone.
  \item Slower movements in shared workspaces.
\end{itemize}

\begin{bigidea}
  Cobots do not replace humans \textemdash{} they become teammates.
  The cobot handles heavy or boring tasks while the human does
  the thinking, checking, and problem-solving.
\end{bigidea}

\sectionrule

% --- Warehouse robots ---
\section*{Warehouse Robots}

When you order something online, a warehouse robot may help find it.
Fleets of small orange robots at companies like Amazon carry entire
shelving units to human packers, saving them thousands of steps a day.

\begin{picturethis}
  Imagine a library where the bookshelves come to \emph{you} instead
  of you walking to the shelves. That is what warehouse robots do
  with products.
\end{picturethis}

\sectionrule


% --- How an assembly line works ---
\section*{How an Assembly Line Works}

\begin{quickmap}
  \textbf{A car on the assembly line:}
  \begin{enumerate}
    \item Robot welder joins the metal body panels.
    \item Robot arm bolts doors into place.
    \item Painting robot sprays three coats of colour.
    \item Camera robot inspects for scratches or dents.
    \item Human worker adds seats, wiring, and trim.
    \item Robot lifts the engine into position.
    \item Quality robot runs a final check.
    \item Finished car rolls off the line!
  \end{enumerate}
\end{quickmap}

\sectionrule

% --- Diagram ---
\begin{diagram}[A simplified assembly line: products move past robot stations.]
  \begin{tikzpicture}[>=Stealth, node distance=2.2cm]
    % Conveyor belt
    \draw[MetalGray,line width=3pt] (-1,0) -- (11,0);
    \draw[MetalGray,line width=1pt,dashed] (-1,-0.15) -- (11,-0.15);
    \node[font=\scriptsize\itshape,below] at (5,-0.3) {Conveyor belt};
    % Stations
    \node[draw=RoboGreen!80!black,fill=RoboGreen!15,rounded corners=4pt,
          minimum width=1.8cm,minimum height=1cm,font=\scriptsize\bfseries,
          above=0.6cm] at (1,0) {Weld};
    \node[draw=WarmOrange,fill=WarmOrange!15,rounded corners=4pt,
          minimum width=1.8cm,minimum height=1cm,font=\scriptsize\bfseries,
          above=0.6cm] at (4,0) {Paint};
    \node[draw=SteelBlue,fill=SteelBlue!15,rounded corners=4pt,
          minimum width=1.8cm,minimum height=1cm,font=\scriptsize\bfseries,
          above=0.6cm] at (7,0) {Inspect};
    \node[draw=SunYellow!80!black,fill=SunYellow!15,rounded corners=4pt,
          minimum width=1.8cm,minimum height=1cm,font=\scriptsize\bfseries,
          above=0.6cm] at (10,0) {Pack};
    % Arrows on belt
    \draw[->,line width=1.2pt,RoboGreen!60!black] (2.2,0.08) -- (2.8,0.08);
    \draw[->,line width=1.2pt,WarmOrange!80!black] (5.2,0.08) -- (5.8,0.08);
    \draw[->,line width=1.2pt,SteelBlue!80!black] (8.2,0.08) -- (8.8,0.08);
  \end{tikzpicture}
\end{diagram}

\sectionrule

% --- Imagine If ---
\begin{picturethis}
  Imagine if every time you made a sandwich, you had to walk to a
  different room for bread, then another for butter, then another for
  cheese. An assembly line brings everything to you in order ---
  that is why factories can build a car in under a minute.
\end{picturethis}

\sectionrule

% --- Try It ---
\begin{tryit}
  Set up a mini assembly line at home.
  \begin{enumerate}
    \item Gather ten identical objects (LEGO bricks, cards, or coins).
    \item Time yourself placing them into a box one by one.
    \item Now ask a friend to pass them to you (like a cobot helping).
    \item Is the teamwork version faster? By how much?
  \end{enumerate}
\end{tryit}

\sectionrule


\begin{quickcheck}
  \begin{enumerate}
    \item What is the difference between a robot and a cobot?
    \item Name three jobs that factory robots do on an assembly line.
    \item How do warehouse robots save workers thousands of steps?
  \end{enumerate}
\end{quickcheck}

\sectionrule
% --- New Words ---
\begin{newwords}
  \begin{description}[labelwidth=3.8cm, leftmargin=4.0cm]
    \item[Assembly line] A production system where products move past stations, each adding a part.
    \item[Welding] Joining two metal pieces by melting their edges together.
    \item[Cobot] A collaborative robot designed to work safely alongside humans.
    \item[Pick and place] A robot task: grab an item from one spot and put it in another.
    \item[Conveyor belt] A moving belt that carries items through a factory.
    \item[Palletising] Stacking finished goods onto a pallet for transport.
  \end{description}
\end{newwords}
